\documentclass[a4paper]{moderncv}

% 风格
\moderncvcolor{cerulean}
\moderncvstyle{contemporary}

% 编码（推荐用 xelatex 编译）
\usepackage{xeCJK}
% \setCJKmainfont{PingFang SC}

\usepackage[scale=0.85]{geometry}

% 个人信息
\name{居}{振梁}
\title{青少年科技辅导员｜课程体系设计师}
\phone[mobile]{+86-186-0123-9387}
\email{juzhenliang@hotmail.com}
\extrainfo{微信：shadowargrey｜公众号：plteenfun}

\begin{document}

\makecvtitle

%--------------------------------------------------
\section{个人简介}

软件工程师出身的\texttt{STEM}教育者，长期从事跨学科课程设计与一线教学。
课程体系以 HtDP 为内核，以信息学奥赛为起点，
强调数学理论与工程实践并重，
并通过自研教学工具支持课程实施与持续优化。
教学经验覆盖2-11年级的多种教学场景，
能有效控制学生认知负荷。

%--------------------------------------------------
\section{教学理念}

在 AI 降低编程门槛的背景下，
教学应以结构理解为核心，
以自顶向下的建模顺序为认知路径，
以真实问题为应用场景，
最终帮助学生建立清晰、可迁移的知识体系。

%--------------------------------------------------
\section{教学与课程经验}

\cventry{2024--2026}{独立教师}{义乌市青舞艺术培训部有限公司}{金华}{}{
\begin{itemize}
\item 构建信息学奥赛编程教育体系，教授4-7年级 CSP/J
\item 优化跨学科编程教育体系，能在2-3年级落地，学生兴致勃勃
\item 与义乌市福田第二小学合作，承担三年级信息技术教学及Python社团课教学
\item 实验性开设高中数学、高中物理、高中技术暑假班，规模不大但客户反馈很好
\end{itemize}
}

\cventry{2021--2024}{教务老师}{镇江市洛克白培训中心有限公司}{镇江}{}{
\begin{itemize}
\item 负责江苏省实验操作大赛第一届和第二届的组织和学生训练工作，金牌、银牌数量总计5个
\item 与镇江培文实验学校合作，3-6年级义务教育科学主教老师，成绩由30-40分提升到70-80分
\item 设计并实施\texttt{STEM}课程体系，覆盖小实验、实验探究及综合项目制课程
\item 构建跨学科编程教育体系，覆盖项目制课程，并衔接蓝桥杯和信息素养大赛
\item 研发教学用游戏引擎及外围工具链，系统性降低了小学生对\texttt{C++}的恐惧
\end{itemize}
}

%--------------------------------------------------
\section{科研与项目经历}

\cventry{2017--2020}{软件工程师}{镇江明润信息科技有限公司}{镇江}{}{
\begin{itemize}
\item 与一线操作人员及领域专家协作，将碎片化经验和工程知识落地成软件
\item 从零构建基于 UWP 平台的疏浚控制系统，以高可靠性替代原有 Win32 系统
\item 开发自适应图形界面框架，将原始传感器数据转化为结构化可视信息（如挖掘臂、GPS 系统等）
\item 参与船载环境连续 3 个月实地测试，在复杂工况下验证系统稳定性，并基于反馈持续迭代优化
\item 总结并形成以“清晰优先于形式”为核心的设计原则，为后续教学导向的软件系统提供方法论基础
\end{itemize}
}

\cventry{2014--Present}{开源贡献者}{}{远程}{}{
\begin{itemize}
\item 研发\texttt{Typed Racket}语义可视化\texttt{DSL}生态，用于表达程序结构与抽象关系，持续演进超过10年
\item 基于函数式编程范式，系统化设计计算机科学教学路径的本土化
\item 将教学引擎架构写成学位论文(查重率 7.73\%，AIGC 1.5\%)
\end{itemize}
}

\cventry{2008--2010}{常熟理工学院——嵌入式网络实验室}{科研助理}{常熟}{}{
\begin{itemize}
\item 协助学生搭建嵌入式开发环境
\item 移植 Linux 内核至目标开发板
\item 开发辅助工具支持学生项目
\item 评审并验证算法与数学模型实现
\end{itemize}
}

%--------------------------------------------------
\section{教学案例}

\cvitem{小学案例}{
\textbf{竞赛路径与学习动机的冲突}
\begin{itemize}
\item \textbf{问题}：学生从\texttt{Scratch/Python}转向\texttt{C++}竞赛训练，学习兴趣显著下降
\item \textbf{设计}：将学习路径拆分为“结构理解 → 半结构问题 → 竞赛问题”，引入过渡层
\item \textbf{实施}：通过自研教学引擎与项目驱动任务，逐步建立抽象能力与问题建模能力
\item \textbf{结果}：在小规模教学实践中能显著提升学习持续性，降低路径切换带来的流失风险
\item \textbf{例证}：五年级男生从抗拒\texttt{C++}到独立完成打靶游戏仅需16次课
\end{itemize}
}

\cvitem{初中案例}{
\textbf{以数学建模为核心的编程教学优化}
\begin{itemize}
\item \textbf{问题}：传统语法驱动路径导致初学者抓不住重点、成长极其缓慢
\item \textbf{设计}：在学习早期就开始培养数学建模意识，让数学和编程的关系名副其实
\item \textbf{实施}：通过精心设计的学生手册引导学生纸笔推演，最终将推演过程翻译成\texttt{C++}代码
\item \textbf{结果}：学生从依赖背模版逐步过渡到独立完成建模任务
\item \textbf{例证}：6年级学生能理解，但下课就忘；7年级学生经教师引导后可迅速上手；8年级学生可仅通过阅读手册提问自己完成建模
\end{itemize}
}

\cvitem{方法论案例}{
\textbf{学习能力评估模型的构建与应用}
\begin{itemize}
\item \textbf{问题}：传统教学评价侧重分数，难以反映学生真实能力结构
\item \textbf{设计}：构建覆盖知识迁移、表达能力、逻辑思维与项目能力的多维评估模型
\item \textbf{实施}：在日常教学与项目实践中嵌入评价维度，用于指导教学决策与个体反馈
\item \textbf{结果}：提升教学针对性，使学生能力发展路径更加清晰可控
\end{itemize}
}

%--------------------------------------------------
\section{培训与证书}

\cventry{2022}{青少年科技辅导员认证}{}{}{}{}

%--------------------------------------------------
\section{教育背景}

\cventry{2020--2025}{南京航空航天大学}{计算机科学与技术}{本科}{南京}{}
\cventry{2006--2009}{常熟理工学院}{J2EE企业级应用开发}{大专}{常熟}{}

%--------------------------------------------------
\section{合作方向}

\cvitem{}{
可与学校及教育机构合作开展以下项目（短期、长期或工作坊）：
\begin{itemize}
\item 跨学科编程融合课程建设
\item 项目制\texttt{STEM}课程开发、优化、实施
\item 编程比赛、信息学奥赛课程
\end{itemize}
}

\end{document}
